% from ega2-1-fr.tex
\oldpage[II]{87}

un entier $k_0>0$ tel que $\mathcal{L}^{\otimes k_0}$ soit très ample
relativement à $q$. D'autre part, en vertu de
(\hyperref[II.4.5.5-fr]{4.5.5}), il existe un entier $m_0$ tel que, pour
$m\geq m_0$, $\mathcal{L}^{\otimes m}$ soit engendré par ses sections
au-dessus de $X$. Posons $n_0=k_0+m_0$~; si $n\geq n_0$, on peut écrire
$n=k_0+m$ avec $m\geq m_0$, d'où
$\mathcal{L}^{\otimes n}
=\mathcal{L}^{\otimes k_0}\otimes\mathcal{L}^{\otimes m}$. Comme
$\mathcal{L}^{\otimes m}$ est engendré par ses sections au-dessus de $X$, il
résulte du critère (\hyperref[II.4.4.8-fr]{4.4.8}) et de
(\hyperref[II.3.4.7-fr]{3.4.7}) que $\mathcal{L}^{\otimes n}$ est très ample
relativement à $q$. Enfin, il est clair que e) entraîne e'), et e') implique
que $\mathcal{L}$ est ample en vertu de (i) et de
(\hyperref[II.4.5.6-fr]{4.5.6, (i)}).

\begin{env}[4.5.10.1]
\label{II.4.5.10.1-fr}
\emph{Démonstration du lemme (\hyperref[II.4.4.10.1-fr]{4.4.10.1}).} Posons
$\mathcal{E}(n)=\mathcal{E}\otimes\mathcal{K}^{\otimes n}$~; comme $g$ est
séparé (\hyperref[II.4.4.2-fr]{4.4.2}), le raisonnement de
(\hyperref[II.3.4.8-fr]{3.4.8}) s'applique et ramène à voir que
l'homomorphisme canonique
$g^*(g_*(\mathcal{E}(n)))\to\mathcal{E}(n)$ est surjectif pour $n$ assez
grand. En outre, comme $Z$ est quasi-compact, le raisonnement de
(\hyperref[II.3.4.6-fr]{3.4.6}) ramène à prouver la proposition dans le cas
où $Z$ est affine. Or, $\mathcal{K}$ est alors ample en vertu de
(\hyperref[II.4.5.10-fr]{4.5.10, (i)}), et la conclusion résulte de
(\hyperref[II.4.5.5-fr]{4.5.5, d')}).
\end{env}
\end{proof}

\begin{corollary}[4.5.11]
\label{II.4.5.11-fr}
Soient $Y$ un schéma affine, $q:X\to Y$ un morphisme séparé de type fini,
$\mathcal{L}$ un $\mathcal{O}_X$-Module ample, $\mathcal{L}'$ un
$\mathcal{O}_X$-Module inversible. Il existe un entier $n_0$ tel que, pour
$n\geq n_0$, $\mathcal{L}^{\otimes n}\otimes\mathcal{L}'$ soit très ample
relativement à $q$.
\end{corollary}

\begin{proof}
En effet, il existe $m_0$ tel que pour $m\geq m_0$,
$\mathcal{L}^{\otimes m}\otimes\mathcal{L}'$ soit engendré par ses sections
au-dessus de $X$ (\hyperref[II.4.5.8-fr]{4.5.8})~; d'autre part, il existe
$k_0$ tel que $\mathcal{L}^{\otimes k}$ soit très ample relativement à $q$
pour $k\geq k_0$. Par suite
$\mathcal{L}^{\otimes(k+m_0)}\otimes\mathcal{L}'$ est très ample pour
$k\geq k_0$ ((\hyperref[II.4.4.8-fr]{4.4.8}) et
(\hyperref[II.3.4.7-fr]{3.4.7})).
\end{proof}

\begin{remark}[4.5.12]
\label{II.4.5.12-fr}
On ignore si l'hypothèse qu'un $\mathcal{O}_X$-Module $\mathcal{L}$ est tel
que $\mathcal{L}^{\otimes n}$ soit très ample (relativement à $q$) entraîne
la même conclusion pour $\mathcal{L}^{\otimes(n+1)}$.
\end{remark}

\begin{proposition}[4.5.13]
\label{II.4.5.13-fr}
Soient $X$ un préschéma quasi-compact, $Z$ un sous-préschéma fermé de $X$
défini par un faisceau quasi-cohérent nilpotent $\mathcal{J}$ d'idéaux de
$\mathcal{O}_X$, $j$ l'injection canonique $Z\to X$. Pour qu'un
$\mathcal{O}_X$-Module inversible $\mathcal{L}$ soit ample, il faut et il
suffit que $\mathcal{L}'=j^*(\mathcal{L})$ soit un
$\mathcal{O}_Z$-Module ample.
\end{proposition}

\begin{proof}
La condition est nécessaire. En effet, pour toute section $f$ de
$\mathcal{L}^{\otimes n}$ au-dessus de $X$, soit $f'$ son image canonique
$f\otimes1$, section de
$\mathcal{L}'^{\otimes n}=\mathcal{L}^{\otimes n}
\otimes_{\mathcal{O}_X}(\mathcal{O}_X/\mathcal{J})$ au-dessus de l'espace
$Z$ (qui est identique à $X$)~; il est clair que $X_f=Z_{f'}$, donc le
critère a) de (\hyperref[II.4.5.2-fr]{4.5.2}) montre que $\mathcal{L}'$ est
ample.

Pour voir que la condition est suffisante, notons d'abord qu'on peut se
ramener au cas où $\mathcal{J}^2=0$, en considérant la suite (finie) des
préschémas
$X_k=(X,\mathcal{O}_X/\mathcal{J}^{k+1})$ dont chacun est sous-préschéma
fermé du suivant, défini par un faisceau d'idéaux de carré nul. D'ailleurs
$X$ est un schéma puisque $X_{\mathrm{red}}$ l'est par hypothèse
((\hyperref[II.4.5.3-fr]{4.5.3}) et
(\hyperref[I.5.5.1-fr]{I, 5.5.1})). Le critère a) de
(\hyperref[II.4.5.2-fr]{4.5.2}) montre qu'il suffira de prouver le

\begin{lemma}[4.5.13.1]
\label{II.4.5.13.1-fr}
Sous les hypothèses de (\hyperref[II.4.5.13-fr]{4.5.13}), supposons de plus
$\mathcal{J}$ de carré nul~; $\mathcal{L}$ étant un
$\mathcal{O}_X$-Module inversible, soit $g$ une section de
$\mathcal{L}'^{\otimes n}$ au-dessus de $Z$ telle que $Z_g$ soit affine.
Alors il existe un entier $m>0$ tel que $g^{\otimes m}$ soit l'image
canonique d'une section $f$ de $\mathcal{L}^{\otimes nm}$ au-dessus de $X$.
\end{lemma}

On a la suite exacte de $\mathcal{O}_X$-Modules
\[
  0\longrightarrow\mathcal{J}(n)\longrightarrow
  \mathcal{O}_X(n)=\mathcal{L}^{\otimes n}\longrightarrow
  \mathcal{O}_Z(n)=\mathcal{L}'^{\otimes n}\longrightarrow0
\]

